\section{Introduction}\label{introduction}

\begin{EssayColumns}

I am tired of watching capable people wait while the organizations
employing them demand more productivity. The work is understood, the
capability is available, and progress still disappears into permissions,
meetings and negotiations over ownership. Companies pay for the people,
the technology and the delay, then ask why AI has not improved the
result.

I began developing \href{https://c89.systems/10x.pdf}{\emph{The
Coordination Shift}} in September 2025 because the economics of my own
work had changed \citep{blomgren2025coordination}. A capable engineer
working with AI could carry a substantially larger and more coherent
body of work than the organizational models around us assumed. More than
a year later, the obstruction is still treated as normal. When
investigating, building and verifying a solution takes less time than
arranging the meeting about whether to begin, the delay is a management
decision.

\NextEssayColumn

Torsten Slok's account of the missing financial payoff and the survey by
Justine Hastings, Chad Syverson and Seth Zimmerman bring that
frustration into focus. Employees already using AI are still spending
much of their working lives making the organization function
\citep{slok2026payoff} \citep{hastings2026deployment}. Leadership cannot
keep demanding the returns of a new production capability while
preserving every dependency of the old one. The response must reach into
who owns decisions, how information moves and what management is
rewarded for maintaining.

\end{EssayColumns}

\EssayArt{opening-illustration.png}

\clearpage

\section{The missing return has an organizational
address}\label{the-missing-return-has-an-organizational-address}

\begin{EssayColumns}

In his October 24 assessment, Apollo chief economist Torsten Slok
contrasts rising revenue per worker and profit margins among the
Magnificent Seven with broadly flat results across the other 493
companies in the S\&P 500. His question is whether cheaper intelligence
will become higher growth and profitability across the wider corporate
economy. He also points to judgment, context, accountability and trust
as obstacles to translating technical capability into business results
\citep{slok2026payoff}. Those are precisely the conditions an
organization has to arrange around productive work.

Three days earlier, Hastings, Syverson and Zimmerman published \emph{The
Deployment Failure Holding Back Your AI Investments} in \emph{Harvard
Business Review}. They surveyed 300 individual contributors and
frontline managers at US companies with annual revenue exceeding \$500
million. Respondents reported spending 37\% of their working time on
coordination, compared with 44\% on substantive work and 18\% on
administration. Every respondent used AI. Yet 76\% wanted to spend less
time coordinating, while 85\% wanted more time for substantive work
\citep{hastings2026deployment}.

\NextEssayColumn

Their definition of coordination is concrete: assembling and cleaning
information from different systems, reconciling conflicting numbers,
refreshing recurring reports and aligning work across functions. The
tools had reached the employees, but much of the effort required to make
the enterprise intelligible to itself remained. An employee could
summarize a document or draft an analysis while still reconstructing the
underlying business situation by hand each time a decision was needed.

The authors distinguish that labour from a second economic cost:
decisions that improve one department's results can damage the
enterprise as a whole. Their product-portfolio example makes the
mechanism clear. Finance sees a product variant that still earns a
profit. Manufacturing sees costly production changeovers. Planning
carries inventory against unpredictable demand, while marketing knows
that many customers would accept a similar product if the variant
disappeared. Each function holds a useful part of the answer. Assessing
the combined effect requires work across all four, while keeping the
product requires no such agreement. The organization repeatedly
authorizes continuity by default \citep{hastings2026deployment}.

\end{EssayColumns}

\SurveyFigure

\clearpage

\begin{EssayColumns}

That default is expensive even when every report is accurate. The
company pays for repeated information assembly, for the time a decision
waits, and for the consequences of decisions made with only part of the
picture. Some losses never appear as a separate expense: an improvement
nobody could authorize, a customer problem nobody owned, or a product
whose full cost was never brought into the decision. Faster
report-writing can leave all three untouched.

Hastings also describes the effect of making recurring analysis a shared
system capability at Amazon. Workforce questions that had required
months of data assembly and statistical work could be answered in days,
leaving scientists and analysts more capacity for difficult cases.

\NextEssayColumn

The article's recommendation follows through to responsibilities and
incentives: leadership must fund work across functions, assign the
decisions that become possible and judge the return across the firm
\citep{hastings2026deployment}. Reducing coordination changes what
people can do; management has to change what they are authorized and
expected to do with it.

Slok's observation and the Hastings analysis therefore belong in the
same discussion. One asks where the financial return has gone; the other
directs attention to the system through which that return must be
realized. My conclusion is that companies should examine the
organization consuming their new capability before concluding that the
capability itself has little value. An unchanged decision system is a
poor instrument for testing the potential of a changed production
system.

\end{EssayColumns}

\PortfolioFigure

\clearpage

\section{The unit of work has
changed}\label{the-unit-of-work-has-changed}

\begin{EssayColumns}

The distinctive claim of the Coordination Shift is that AI changes the
scale of work a person or small group can own. When analysis,
implementation, verification and revision can remain within a coherent
human--AI unit, some of the handoffs previously needed to combine those
activities disappear. The gain extends beyond faster execution of an
individual task to the possibility of completing work that previously
depended on several separate contributions.

Experimental evidence gives that proposition a concrete basis. In
\emph{The Cybernetic Teammate}, Dell'Acqua and colleagues studied 791
Procter \& Gamble professionals working on product-innovation
challenges. Individuals using AI matched the performance of two-person
teams without it. AI-assisted participants also produced more balanced
technical and commercial proposals, crossing distinctions that otherwise
followed professional specialization. Some benefits obtained through
collaboration could be reproduced within an individual's augmented work
\citep{dellacqua2026teammate}.

The organizational boundary appears just as clearly in Dillon and
colleagues' \emph{Shifting Work Patterns with Generative AI}. Their
six-month randomized experiment covered approximately 6,000 knowledge
workers. Workers who used the tool spent around three fewer hours per
week on email, while time in meetings did not change significantly.
Activities individuals could alter independently changed more readily
than activities requiring coordination with others
\citep{dillon2025workpatterns}.

\NextEssayColumn

Those findings describe an expanding capacity to act inside an
arrangement that is harder to change. An employee can finish the
analysis sooner without gaining authority to implement its conclusion.
An engineer can build a solution without changing the process for
obtaining data, agreeing ownership or releasing it. The organization
receives more completed contributions, but the contributions still wait
to become a completed outcome. Where review and decision capacity remain
fixed, increasing the arrival rate can enlarge the queue.

The industrial division of labour rests partly on a trade:
specialization improves performance, while coordination pays for
bringing the specialists' contributions together. As AI expands the
useful range of individual expertise, that trade must be recalculated. A
boundary established because several people were once required for a
piece of work does not remain economically justified merely because it
has acquired a manager, a budget and a process description.

An ordinary AI programme can miss this opportunity by beginning with
existing jobs and asking how to accelerate their tasks. The more
consequential question is which tasks, dependencies and divisions are
still necessary when those capabilities are combined differently.
Helping six departments communicate faster can be valuable. Enabling a
coherent unit to complete the work without four of those exchanges
changes the operating model itself.

\end{EssayColumns}

\UnitComic

\clearpage

\section{Why the company keeps
paying}\label{why-the-company-keeps-paying}

\PrepareGutterQuote{14}{78mm}{The institution can afford\\the delay while the useful\\work remains undone.}

\begin{EssayColumns}\GutterQuote

A company can cover its costs, retain its customers and satisfy its
owners while operating far below its available productive capacity.
Contracts, switching costs and established relationships can protect an
incumbent whose competitors carry similar overhead. A more capable
entrant must still make its offer accessible and credible to buyers.
Technical superiority does not abolish the customer's procurement
process, so the pressure to reorganize may arrive much later than the
ability to do so.

Persistent inefficiency has an empirical history. In \emph{Does
Management Matter? Evidence from India}, Bloom and colleagues found that
a randomized management intervention increased productivity in textile
plants by 17\% during the first year. Informational barriers had impeded
adoption of profitable practices, while limits on managerial time
constrained expansion and reallocation sufficiently that competition had
not forced badly managed firms to exit \citep{bloom2013management}.
Available improvement and market survival were compatible with continued
underperformance.

Inside a large enterprise, the old arrangement has another advantage: it
is already authorized. Established headcount, reporting routines and
approval chains often require less explicit justification than changing
them. A new AI subscription has a visible price and an owner expected to
demonstrate its return. The accumulated cost of people waiting on one
another is dispersed through existing budgets. Management can scrutinize
the new expense while treating the larger inherited expense as an
unavoidable condition of doing business.

\NextGutterColumn

The incentives can reinforce the arrangement without anyone deliberately
choosing waste. A department protects its budget, a manager protects a
commitment, and a specialist avoids accepting responsibility beyond
their mandate. A recurring fault may remain unfixed because the team
that would pay for correcting it is not the team absorbing its operating
cost. Each decision can be defensible within its own boundary. The
enterprise pays for the combined result, and exhortations to collaborate
leave the distribution of authority untouched. The coordination trap
closes when removing these dependencies requires agreement from the same
divided functions, each judged against the local commitments that made
the dependencies expensive in the first place.

Syverson's earlier work supplies a useful connection. In \emph{The
Productivity J-Curve}, with Erik Brynjolfsson and Daniel Rock, he
examined the complementary investments needed to realize value from
general-purpose technologies, including intangible investments poorly
reflected in conventional measures \citep{brynjolfsson2021jcurve}. That
is an argument for doing the organizational work. The passage of time
does not assign a decision owner, connect inaccessible information or
remove a redundant approval. A company expecting the benefits of
reorganization should be able to identify what it is actually
reorganizing.

Companies may continue paying the coordination cost for as long as their
economics permit it. That makes adaptation a leadership decision well
before it becomes a condition of survival. For employees who can see the
unused possibilities, an acceptable operating margin does little to make
the experience less absurd. The institution can afford the delay while
the useful work remains undone.

\wrapstuffclear\end{EssayColumns}

\clearpage

\section{Give the whole outcome an
owner}\label{give-the-whole-outcome-an-owner}

\PrepareGutterQuote{14}{82mm}{A unit that must seek fresh\\agreement across departments\\for every meaningful action has\\responsibility without control.}

\begin{EssayColumns}\GutterQuote

The response should begin where Hastings and colleagues locate the
problem: at the boundaries between information, decisions and
incentives. A complete business judgment needs both a dependable picture
of its consequences and a path to action. Giving each department a
faster way to produce its contribution leaves that path fragmented.
Reorganizing around a coherent outcome makes the fragmentation itself
something management must remove.

There is an established foundation for this approach. Michael Hammer's
\emph{Reengineering Work: Don't Automate, Obliterate} challenged
companies using information technology merely to accelerate outdated
processes \citep{hammer1990reengineering}. The Coordination Shift brings
a specific new reason to revisit that challenge: the human--AI
production unit can internalize work that an earlier organization had to
distribute. The opportunity concerns both the process and the division
of responsibility behind it.

My original analysis examined organizational models for their
compatibility with this change. Team Topologies, platform organizations,
networked or cellular structures and autonomous micro-enterprises offer
useful foundations because they support local ownership, explicit
boundaries and dependable shared capabilities
\citep{blomgren2025coordination}. Their practical contribution is to
reduce recurring negotiation. The appropriate arrangement follows the
outcome and required competence, including continuity and operational
coverage.

\NextGutterColumn

Team Topologies distinguishes time-bounded collaboration, service
provision and facilitation \citep{teamtopologiesconcepts}. People may
need to work intensively together to establish a capability; consuming
the established capability should not require repeating that
collaboration indefinitely. A shared data service, deployment platform
or assurance capability earns its place by making other units more
independent. If ordinary use still requires a succession of tickets and
negotiations, the dependency has merely acquired a technical interface.

The precedents are tangible. In REA Group's account of its
financial-services transformation, teams became smaller, domain
ownership clearer and decisions more decentralized. Its leaders reported
doubled flow efficiency and improved first-time-right delivery
\citep{devine2026rea}. They changed the conditions under which work
could proceed. Management should scrutinize those conditions with the
same seriousness it applies to the performance of the people working
inside them.

Return to the product-range decision. A unit accountable for the
economic health of one product family needs access to commercial,
production, inventory and service information. It also needs a mandate
to change defined aspects of the offering or operation within agreed
limits. Domain expertise establishes what the information means;
technical capability makes the relevant information and actions
reliable; AI expands the analysis and execution they can perform
together. Responsibility continues through the resulting change and its
observed business effect.

\wrapstuffclear\end{EssayColumns}

\clearpage

\begin{EssayColumns}

The mandate must move real decisions. Which stock policies may the unit
adjust? Which variants may it retire within its authority? What customer
obligations and financial limits govern the work? Who decides an
exception, and how quickly? These questions turn autonomy into an
operating arrangement. A unit that must seek fresh agreement across
departments for every meaningful action has responsibility without
control.

Leadership remains responsible for enterprise direction, investment and
trade-offs between outcomes. It must prevent units from improving their
own figures by exporting costs or risks elsewhere. Finance, security,
architecture and other specialist functions help establish common
constraints and provide expertise when it is needed. Their contribution
becomes more valuable when it improves the conditions under which many
decisions can be made, instead of repeatedly granting permission for
substantially similar decisions.

\NextEssayColumn

Begin with one bounded outcome and a sponsor who can change the relevant
mandates and budgets. Establish the current elapsed time, total effort,
decision delays and business result. Give the unit the access, authority
and shared services it needs, then remove the handoffs that the new
arrangement replaces. Keeping every old approval while adding an
empowered team simply creates a more expensive version of the original
organization.

Evaluate the complete result: customer service, contribution, quality,
recurrence of problems and full operating cost, including AI use and
verification. Decide how released capacity will be used. It might
support additional business, stronger reliability, lower expenditure,
learning or better working conditions. Capacity becomes value through a
deliberate use; a dashboard reporting hours saved cannot make that
decision for management.

The transition also needs continuity. Shared records, documented
decisions, executable checks, practiced recovery and apprenticeship make
capability transferable. As successful arrangements spread, retire
superseded reporting obligations and reward managers for removing
recurring work. The test is whether the new way of operating becomes
ordinary, with the authority and support to continue after the first
sponsor turns their attention elsewhere.

\end{EssayColumns}

\OutcomeFigure

\clearpage

\section{Safety must be able to act}\label{safety-must-be-able-to-act}

\begin{EssayColumns}

The same coordination problem has consequences beyond productivity.
Safety work requires observations to reach someone who understands them
and can cause an effective response. If evidence, judgment and authority
sit behind separate queues, an organization can devote considerable
effort to safety while failing to control a developing hazard. A
documented concern remains unresolved until the necessary action has
taken effect.

Leveson and Thomas's \emph{STPA Handbook} treats the timing of control
actions as part of safety. An action can be unsafe because it arrives
too late, in the wrong sequence or for the wrong duration
\citep{leveson2018stpa}. The implication for organizational design is
direct: the path between recognizing danger and exercising control
belongs to the safety system. Its delay cannot be dismissed as merely
administrative.

Software-delivery research offers a concrete example. DORA's findings on
change approval, drawing on its 2019 research, associate heavyweight
external approvals with poorer delivery performance and find no evidence
that more formal external review produces lower change-failure rates.
Its guidance places peer review close to delivery and supports it with
automated testing and monitoring \citep{dorachangeapproval}. The value
of a control depends on whether it prevents, detects or contains a
harmful change.

\NextEssayColumn

A useful governance architecture establishes boundaries before routine
action is needed, enforces them during operation and preserves evidence
for independent scrutiny. Permissions, verification, monitoring,
containment and recovery should be available to the people accountable
for the outcome. Human deliberation determines consequential risk
choices and changes to those boundaries. Operational responders need the
authority to stop or correct an activity within them, including when
waiting would increase harm.

Slowing or stopping execution can be the correct control action. The
dangerous delay is the inability to decide and enforce that action while
exposure continues. An organization whose safety response depends on a
chain of unavailable decision-makers has built latency into its
protection. As more consequential activity becomes automated, that
weakness can affect employees, customers and the public. A safety
function should be able to show how its controls reduce exposure and
improve intervention, including how quickly a known hazard can be
brought under control.

\end{EssayColumns}

\SafetyFigure

\clearpage

\section{Stop financing the
obstruction}\label{stop-financing-the-obstruction}

\PrepareGutterQuote{8}{78mm}{We cannot keep paying\\capable people to wait\\and expect abundance\\to emerge from the queue.}

\begin{EssayColumns}\GutterQuote

The two articles that prompted this essay belong together. Slok
describes a corporate economy in which much of the anticipated payoff
remains difficult to see. Hastings, Syverson and Zimmerman locate a
substantial burden in coordination and fragmented decisions inside firms
already using AI. The Coordination Shift connects those observations to
a change in what individuals and small units can accomplish. New
capacity is entering organizations still arranged around older
constraints.

The response is available in practical organizational design: coherent
outcome ownership, authority where the relevant knowledge resides,
shared services that remove dependencies, and governance capable of
acting. These changes require leadership because they alter budgets,
responsibilities and the justification for existing processes.

\NextGutterColumn

They will not emerge automatically from better prompts or another model
release. Companies have to choose whether increased capability changes
the institution or is absorbed by it.

My impatience comes from the useful work being left undone while we
discuss how transformative the technology might become. Better services,
affordable solutions and less exhausting work are concrete
possibilities. We are discussing an age of abundance while tolerating
organizations that make useful action scarce. Productive machines need
institutions prepared to let that productivity reach the world. We
cannot keep paying capable people to wait and expect abundance to emerge
from the queue.

\wrapstuffclear\end{EssayColumns}

\EssayArt{closing-illustration.png}
